\section{Theoretical Framework}

\utsett{chapter opening paragraph 2}


%-----------------------------------------------------------------------
\subsection{Energy system Flexibility: Concepts and Frameworks}
\utsett{sub-section opening paragraph 2.1}

\subsubsection{Definitions and Frameworks}

No single definition of energy flexibility commands agreement across the heating and power-system literature, and the variation is consequential because each definition draws a different boundary around what counts as flexibility. Two recent reviews document the fragmentation directly: A SINTEF synthesis grades $14$ prior definitions against three criteria, find none that satisfies all, and proposes a unifying definition to close the gap \cite{degefa_comprehensive_2021}. 
\utfyll{Write in what that proposition is depending on how it relates to DH}
A Systematic review of $85$ building-flexibility studies finds that roughly $42\%$ of them give no formal definition of flexibility, and that the rest rely on $61$ distinct quantification metrics \cite{li_energy_2021}. Three broad traditions recur. The first treats flexibility as a physical property of the system, grounded in inertia and in the coordinated operation of coupled electricity, heat and gas grids: That flexibility is simply the ability to speed up, delay or reduce the injection or extraction of energy relative to the reference profile, requiring only that the system possess inertia \cite{vandermeulen_controlling_2018}, \cite{lund_smart_2017}, \cite{mathiesen_smart_2015}.
The second defines flexibility operationally, as a response to an external signal: The managed shifting of demand and generation to meet grid and local objectives, quantified through market-tier participation or performance indicators such as the "Flexibility Function" rather than in absolute energy terms \cite{li_ten_2022}, \cite{golmohamadi_integration_2022}. 
The third, sociological, tradition argues that both framings take the timing of demand as given, and so risk stabilizing rather than providing a genuine flexible system \cite{blue_conceptualising_2020}, a caveat this thesis acknowledges while remaining within the engineering and economic frameworks. The traditions converge on flexibility as the temporal shifting of energy injection or extraction enabled by inertia, and diverge mainly on whether it is treated as a system property, a market tradable response, or a control problem. This thesis adopts a definition that sits in the first two traditions, and suits a district heating operator facing the power market. The thesis will operate under the following definition: Energy flexibility is the capacity of an asset or system to modify its rate of production or consumption relative to a reference profile in response to a price- or system signal, within technical and commercial constraints \cite{sneum_barriers_2021}, \cite{degefa_comprehensive_2021}. 
\utfyll{Something along the lines of: If a flexible asset is never activated, due to either non-existent pricing signals, or if market/tariff structures are set up in such a way as to cost more for the asset owner than it earns in providing the flexibility - then the flexibility asset only holds theoretical societal value, which cannot be utilized until regulatory changes are made.}

A second distinction of flexibility assets separates between the physical capacity for flexibility, and its operation organizes how that capacity is described throughout this thesis. Following Vandermeulen in \cite{vandermeulen_controlling_2018}, the \textit{structural} layer is the embedded physical capability that makes deviation from the reference profile possible. In this context, it is the of thermal inertia in a district heating system through the heat carrier in the network pipes, dedicated thermal energy storage such as an accumulator tank, the thermal mass of connected buildings and electric boilers. The \textit{operational} layer is the dispatch decision that activates the physical capacity in response to a signal. These two are utilization layers of the same technology; an electric boiler is structural in its installed capacity, but the choice to up- or down-regulate its usage is operational. Following this, we can state that structural capacity without a operational dispatch signal or utilization yields no flexibility. 

Structural capacity and operational dispatch together constitute a \textit{flexibility resource}, and whether that resource provides any value to the operator depends on the third layer, \textit{enablers}. Degefa \cite{degefa_comprehensive_2021} and Sæle \cite{saele_understanding_2023} classify markets, reserve products and network tariffs as flexibility enablers - the external preconditions required to activate the value of flexibility resources. A district heating operator can hold ample structural flexibility capacity and even dispatch it correctly, yet capture no value if no reserve product, bilateral contract or proper market conditions exists. 
\utsett{AUTOREF: Section 2.3 catalogues the enablers that are missing or adverse in Norway, and AUTOREF section 2.5 formalizes how their absence creates a value gap between the socio-economic benefits and the commercial value of the same physical flexibility}

Operational flexibility is exercised in two modes. The first is load-shifting; advancing or deferring consumption relative to the reference profile without altering the total energy delivered - it repositions the time at which the same amount of consumption occurs. The second, and the principal operational lever in district heating with electrical boilers installed, is directional regulation - the increase or decrease of electrical consumption in response to system conditions. The Nordic flexibility framing of the Flex4RES program defines flexibility in these directional terms; As the ability to regulate the increase or decrease of generation or consumption in response to market signals \cite{skytte_flexible_2019}. 
The two directions map onto opposite system needs. When electricity is in surplus and prices are low, a power-to-heat (P2H) asset up-regulates - increasing electrical consumption to convert surplus power into heat and charging the accumulator tank for later use. When the grid is stressed and prices are high, the same system down-regulates - reducing electric boiler draw and instead discharging stored heat from the tank, which relieves the peak load on the grid while still serving demand.

The power flexibility of a district heating system with a controllable electrical boiler is sentral to this thesis. Golmohamadi distinguishes power flexibility, a response in the system's electricity demand, from heat-side flexibility, and names the electrically operated assets as the source of the greatest direct power-flexibility potential in district heating systems \cite{golmohamadi_integration_2022}. Bloeß and Schill frame the same flexibility as the \textit{flexible} use of electricity for heating, the timing of the conversion adjusted to system conditions rather than electricity simply substituting for fuel \cite{bloes_power--heat_2018}. 

How far that power flexibility can be shifted in time depends on the thermal buffer into which the converted heat can be stored, and later drawn from. Golmohamadi groups three such buffers into two kinds \cite{golmohamadi_integration_2022}:
\begin{itemize}
  \item Passive thermal inertia:
  \begin{enumerate}
    \item The heat carrier in the network pipes (FHC, real-time flexibility/intraday)
    \item The thermal mass of connected buildings (FTI, real-time flexibility/intraday)
  \end{enumerate}
  \item Active thermal inertia:
  \begin{enumerate}
    \item Dedicated thermal storage (FTS): An accumulator tank that scales from day-ahead arbitrage to seasonal shifting depending on its size.
  \end{enumerate}
\end{itemize}
For a district heating operator working from th supply side, the relevant buffers are the network heat carrier (FHC) and the accumulator (FTS). Building thermal inertia (FTI) is a demand-side source outside the scope of this thesis. The buffer's parameter; heating storage capacity, sustainable deviation duration, and the ramp rate at which changes can be applied, determine which products the power-flexibility action can bid into. Two revenue pathways appear: The first pathway is reserve products, procured by the transmission system operator and graded by its activation speed, pay for capacity held ready and energy activated on a system signal. The second being spot-market arbitrage which rewards shifting consumption across the day-ahead price profile, and requires a price spread between charging and discharging windows wide enough to cover the cost of flexible operation. Each pathway faces different eligibility barriers, activates on a different signal and forms a separate component of the commercial value of thermal flexibility.

Just as the definition of flexibility differs, so does the estimation of its value. The term is applied to everything from reduced system operating cost or welfare gain in \cite{mitridati_heat_2020}, to plant-level arbitrage revenue in \cite{javanshir_operation_2022}, to avoided variable renewable energy curtailment and the corresponding uplift in its generation value \cite{kirkerud_power--heat_2017}, as well as building or network performance indicators such as peak-to-average load ratio \cite{li_ten_2022}. These framings apply different system boundaries, count different actors as beneficiaries, and produce results in incomparable units. This thesis adopts a perspective-explicit working definition: Flexibility value is the monetary or welfare benefit attributed to a flexibility action, conditional on the system boundary and the actor whose ledger is being measured. \utsett{AUTOREF: Section 2.5} formalizes this as the commercial interest and socio-economic value perimeters. The following two subsections establish the Norwegian power system and the sector-coupling context within which both perimeters operate.

%-----------------------------------------------------------------------
\subsubsection{Flexibility in Power Systems}

The Norwegian power system is organized into five bidding zones, NO1 through NO5, which clears with their own separate spot prices through NordPool. The zones reflect physical transmission constraints between regions, and is a regulatory instrument for managing congestion - when transfer limits are met, each zone clears its own price, which signals local scarcity or surplus \cite{iea_-_international_energy_agency_norway_2022}. Hydropower supplies close to $90\%$ of the national generation, making it the dominant source in every pricing zone. NO2, the southwest, is the largest producer and runs a substantial surplus, together with the western (NO5) and northern (NO4) who also generates above their local demand. NO1, the southeastern grid including Oslo, is the demand center of the country, and is generally a net importer of power. 
\utsett{concrete figures around import vs export vs production type in the different regions - Electricity consumption and generation by price area (StatBank table 14092 \& 14093 SSB: September 2025, for instance, it generated about 1.6 TWh against 2.1 TWh of consumption, with roughly 94\% of its own output was hydropower, whereas NO2 produced about 3.8 TWh against 2.5 TWh of demand}.

Southern Norway (NO1, NO2, and NO5) are strongly coupled and frequently clears at a single price. The sharper divide runs between the south and the wind-heavy north (NO3 and NO4), not the southern neighbor-zones \cite{iea_-_international_energy_agency_norway_2022}. Because NO1 is closely coupled to NO2, a wind-driven low price surplus in the south propagates much of its market signal to the Oslo zone. The arbitrage value available to a power-to-heat operator in NO1 is therefore governed less by local conditions and more by the southern wind and system coupling to continental Europe, and by congestion in the connector to the northern zones. 
\utsett{Analyze NO1 spot-price and day-ahead price variance and quantify NO1's coupling to NO2+NO5 zones, and the NO3+NO4 zones. Adjust the paragraph acordingly with the nordpool and ssb datasets}
The resulting price-spead distribution sets the ceiling of the commercial arbitrage value available to Oslo-based DH operators before any market-access barriers are taken into account.

Norway's hydropower dominance reshapes the value logic of thermal flexibility relative to fossil-heavy or wind-heavy power systems \cite{NVE_langsiktig_2023}. The combination of dispatchable hydropower and price-responsive industrial demand makes Norway the Nordic region's primary provider of system flexibility \cite{skytte_flexible_2019}. In the terms of \utsett{AUTOREF Section 2.1.1}, a hydropower reservoir is itself a structural flexibility asset at system scale. In a fossil-dominated system, a load that absorbs off-peak electricity displaces marginal thermal generation, reducing fuel combustion and the associated emissions. In Norway, where the marginal unit originates from hydropower, drawing power off-peak rather than at the demand peak conserves the reservoir water for higher-value uses, namely peak-load dispatch, cross border export to higher priced markets, and reserve provision to Statnett \cite{askeland_balancing_2019}, making the primary system-level benefit of sector coupling in Norway freed reservoir capacity rather than saved fossil fuel cost and consumption. Empirically, power-to-heat in the Nordic district heating mix concentrates in the spring, summer and autumn rather than during peak winter demand because the high electricity price during winter hours cause other heating fuels to become cheaper than electric conversion \cite{kirkerud_power--heat_2017}. 

Statnett procures balancing services through a three-tier product hierarchy that cascades through increased time-frames, from automated to manual activation. Frequency Containment Reserves (FCR) are activated continuously and symmetrically in response to real-time frequency deviations. Automatic Frequency Restoration Reserves (aFRR) restore nominal frequency automatically within seconds to minutes following a disturbance. Manual Frequency Restoration Reserves (mFRR) are dispatched manually on a 15-minute window and constitute the product tier most accessible to large thermal demand-side assets such as district heating electric boilers. Until early 2025 the minimum bid size for mFRR participation was $10 MW$ nationally and $5 MW$ in the NO1 and NO3 bidding zones \cite{monie_power--heat_2022}, a threshold that excluded all but the largest aggregated demand-side assets. With the implementation of the Nordic mFRR Energy Activation Market on the 4th of march 2025, the minimum bid was harmonized to $1 MW$ across all Nordic bidding zones \cite{solberg_manedsrapport_nodate} - a change Statnett had signaled in advance as part of its transition to automated balancing \cite{dalen_fleksibilitet_nodate}. The bid-size floor therefore no longer functions as a structural exclusion mechanism for district heating, where a single electric boiler site now clears this threshold. The binding barriers to district heating participation shift accordingly to qualification cost, and operational and metering complexity. The absence of a market framework along with the consumption capacity charge remain the largest barriers to DH flexibility utilization. The remaining barriers are further examined in \utsett{AUTOREF section 2.3.2}, and quantified in \utsett{AUTOREF section 4.3}.


%-----------------------------------------------------------------------
\subsubsection{Sector Coupling and Thermal Systems}

Sector coupling denotes the intentional integration of distinct energy-carrier systems through conversion technologies that let energy flow in parallel across historically separate supply chains \cite{lund_smart_2017}. The operative coupling in this thesis is power-to-heat: the conversion of electrical energy to thermal energy through electric boilers and heat pumps, or substitution of building heating with power from the grid through substituting it with district heating. The electrical. Boilers and heat pumps creates a controllable load on the power system and a corresponding heat source within the district heating network, making the network an active interface between the two carriers rather than a passive heat consumer. For both technologies, the coupling is unidirectional, in that the electricity is consumed to produce heat with no reverse conversion within the asset. Combined heat and power (CHP), including waste incineration plants couples the same two energy carriers, but in the opposite direction - generating electricity alongside heat from fuel. Its flexibility comes from adjusting the fuel consumption or the heat-to-power ratio. In the Oslo system, it serves as a baseload rather than as a P2H flexibility asset. Realizing the value of the P2H coupling depends on policy and market design as much as on the conversion technology. A scoping review of sector-coupling research by Vanhanen and Hanhijärvi \cite{vanhanen_energy_2024} found that policy integration and sector coupling are studied largely in parallel, with district heating being the explicit focus of only one of the $25$ reviewed studies, leaving the coordinated policy conditions for P2H flexibility under-examined.

The structural flexibility of a district heating system, its thermal inertia is what makes this interface relevant at meaningful scale. Network piping volume and dedicated accumulator tanks can absorb or release heat independently of instantaneous consumer demand, giving the operator a thermal buffer across which electrical load can be advanced or deferred. The structural capacity of that buffer, with its volume,temperature range, and charge and discharge rate sets the envelope within which operational load-shifting if feasible. When electricity prices are low, a P2H asset converts surplus power to heat and stores it in that buffer, to be released during high-price or high-demand periods, which can help free up grid capacity and attenuate price spikes and congestion. The paper \cite{mitridati_heat_2020} name this a virtual electricity storage function, in which the district heating system provides energy storage-equivalent services to the power system, by synchronizing its electrical load with system conditions, without a dedicated power-storage asset. They show that sequential heat-then-electricity market clearing is shortsighted, so an operator optimizing within its own heat-market perimeter forgoes the system value that coordinated clearing would realize. The storage-equivalence itself is quantified in \cite{mitridati_power_2018}, which models the district heating pipeline storage as a power-system resource on a six-bus, three node case. The paper reports a $3.51\%$ reduction in total cost and a $1.65\%$ uplift in wind utilization when the virtual storage capability is unlocked. These magnitudes are case-specific, but the results indicate that heat storage-equivalency has merit, and that district heating thermal inertia carries quantifiable system benefit when the market design permits coordinated clearing. The critical qualification is that this value must be accessible to the operator who provides it. This accessibility hinges on an enabler; a reserve market product, flexibility contracts, or market and fee structures which rewards coordinated dispatch. Without an enabling function like that, the value accrues only to the grid operators and consumers, while the operator bears the full cost of flexible operation. This asymmetry between the location of value-creation and the location of cost incidence is the central theoretical mechanism this thesis seeks to quantify and reduce. 

The Norwegian consequence of this coupling value differs from other European nations in one important aspect. In systems with high shares of variable renewable generation (VRE), the primary system-level benefit of P2H is the absorption of surplus wind and solar energy which would otherwise be curtailed, thus raising the achievable penetration of renewable generation \cite{geabermudez_role_2021}. In Norway, where current curtailment is structurally low and the dominating marginal power unit is hydropower reservoirs, the main benefit is instead reservoir-conservation. Surpluss absorption frees water for peak dispatch, export or reserve provision rather than averting VRE curtailment \cite{askeland_balancing_2019}. The coupling mechanism is identical across both contexts, but the value pathway differs. \utsett{AUTOREF: Section 2.4} develops the Norwegian system characteristics that produce these dynamics, and \utsett{AUTOREF: section 2.5} places it within the commercial DH and socio-economic valuation perimeters.


%-----------------------------------------------------------------------
%-----------------------------------------------------------------------
\subsection{District Heating as a Flexibility Resource}

\utsett{Subchapter introduction 2.2}

%-----------------------------------------------------------------------
\subsubsection{European and Nordic Overview}

District heating covers a significant, but unevenly distributed share of building heat demand across Europe. Throughout this thesis, the penetration metric is the share of national buildings heat demand meet by district heating. The most comprehensive cross-country baseline on this metric is close to a decade old, but because district heating penetration is is governed by long-lived network infrastructure, it shifts slowly. According to that study, Denmark leads at roughly $50-65 \%$, with Sweden and Finland following at comparable levels. Germany, the UK and most of southern Europe remain below $15\%$ \cite{werner_international_2017}. A more recent study corroborates the cross-country ranking, while observing that the Danish DH meets roughly two-thirds of residential heating demand by the end of the 2010's \cite{johansen_something_2022}. The sector is undergoing a generational transition towards lower supply temperatures and greater integration of renewable and waste incineration sources. The fourth-generation (4GDH) concept that frames this shift was defined over a decade ago \cite{lund_4th_2014}. A 2024 comprehensive survey extends the trajectory to a fifth generation (5GDH) of near-ambient temperature networks, which is still evolving alongside 4GDH, rather than replacing it, and remains confined to early pilot projects \cite{dang_fifth_2024}. This thesis concerns existing large-scale DH infrastructure, operating at conventional supply temperatures. 4GDH and 5GDH developments are noted as background context, but do not form a part of the analysis. 

Denmark occupies a disproportionately large share of the DH flexibility literature relative to its system scale, reflecting its early large-scale P2H deployment and a long tradition of integrating district heating with wind power. The Danish sector rests on statutory advantages that Norway does not share: Mandatory connection for new development, a non-profit pricing principle, and long-term municipal financing and ownership that aligns operator incentives with broader energy-system objectives \cite{johansen_something_2022}. Greater Copenhagen, where large-scale heat pumps have been integrated into the district heating network, has become a recurring technical benchmark for DH-integrated P2H flexibility \cite{bach_integration_2016}. Many of the barrier typologies and analytical frameworks applied in this chapter were developed in the Danish context, and are transferred here with adjustment for Norwegian structural conditions.

Sweden represents a comparatively under-examined Nordic case. Swedish district heating covered approximately $50\%$ of national building heat demand in 2017, with Stockholm, Gothenburg and Malmö operating large network-integrated systems \cite{werner_international_2017}. The SthlmFlex project demonstrated a local flexibility market architecture within the Stockholm system, initiated jointly by Svenska kraftnät and the regional DSOs Vattenfall Eldistribution and Ellevio to address congestion that occurs during the winter-season \cite{nyman_how_2026}. The pilot operated from winter 2020/2021 onward and has since concluded with a very low market liquidity, with Ellevio stating that short-term local flexibility markets are not a suitable congestion-management tool for Stockholm in the near term future \cite{scrocca_local_2026}. The pilot nevertheless remains the closest Nordic institutional reference pilot for multi-DRO/TSO flexibility coordination. \cite{atabaki_optimising_2026} offer a recent sector-coupling optimization for the Swedish DH sector that is structurally analogous to this thesis's Norwegian question, and \cite{fernqvist_district_2023} provide the only available qualitative study of Swedish DH-flexibility actors - finding that weak institutional incentives and profit-oriented ownership structures limit flexibility provision, despite adequate technical capacity. Sweden's combination of large urban DH systems, growing wind integration and emerging local flexibility market structures makes it the more instructive Nordic comparator for analysis in \utsett{AUTOREF: Section 2.3.5}, wheras Denmark provides the primary technical benchmark throughout \utsett{AUTOREF: Sections 2.2 and 2.3}.

Norway presents a structurally distinct profile within the Nordic group. National district heating penetration is approximately 10 to 12 percent of national building heat demand, significantly lower than Danish or Swedish levels, with the sector concentrated in three cities with three separate operators: Oslo (Hafslund Celsio), Trondheim (Statkraft Varme), and Bergen (BKK Varme) \cite{sandberg_framework_2018}. The paper attributes this gap to historical lock-in rather than to differences in operator profitability: Norway's legacy of abundant, historically cheap electricity and a direct-electric heating culture left it with few water-borne distribution systems, so the country now combines the lowest overall DH penetration in the Nordic region with the highest power-to-heat share of district heating generation \cite{sandberg_framework_2018}. Norwegian DH operators must consequently compete against a low-cost electricity baseline and lack the scale, cooperative ownership structures, and policy support that sustain Danish and Swedish operators. This structural position shapes both the barrier analysis in \utsett{AUTOREF: Section 2.3} and the case-selection rationale in \utsett{AUTOREF: Section 3.1}


%-----------------------------------------------------------------
\subsubsection{Power-to-Heat Technologies}

Power-to-heat (P2H) technologies are the controllable electrical loads that couple the electricity system to the district heating network: electric boilers and large-scale compression heat pumps convert purchased electricity into heat and can therefore be dispatched against electricity price signals and reserve-market activation windows \cite{bloes_power--heat_2018}, \cite{golmohamadi_integration_2022}. A third heat source, combined heat and power (CHP) plants including waste-to-energy (WtE) incineration facilities, couples the two carriers in reverse, generating electricity as a byproduct of heat production rather than consuming electricity \cite{beiron_combined_2020} and is the baseload technology for Oslo's district heating \cite{sandberg_framework_2018}. It is treated here because it sets the merit-order floor around which the P2H assets dispatch, not because it is itself a P2H flexibility asset. The three sources differ in capital intensity, operating cost structure, ramp rate, and the mechanism through which they bear on flexibility provision, distinctions that are load-bearing for the commercial model in \utsett{AUTOREF section 3.3}.

Electrode boilers convert electrical current directly to heat by passing it through a water or steam medium, a process whose ramp rate is limited in practice only by switchgear response times rather than by thermal inertia, which makes the electrode boiler the fastest-responding P2H asset. Capital cost is low among P2H technologies, but operating cost is dominated by the electricity purchase price, making the asset economically attractive only when spot prices are sufficiently low or when reserve capacity payments compensate for standby commitments \cite{bloes_power--heat_2018}.

Large-scale district heating installations are predominantly medium-voltage electrode units in the $3-70 MW$ range that modulate output continuously, in contrast to the lower-voltage resistance boilers and older stepped designs that operate across discrete load steps; both types reach conversion efficiencies of $95-99\%$ \cite{maruf_classification_2022}. In a statement to the online magazine Energiteknikk, Helle Grønli from AFRY points to a value-add on two fronts: Electrode boilers could earn meaningful annual revenue from mFRR capacity payments, while optimizing dispatch against the spot price could deliver sizable operational cost savings, mostly from heat-pump operation \cite{gronli_hvor_2025}. 

Large-scale compression heat pumps extract thermal energy from a low-grade source, such as seawater, groundwater, or wastewater, and upgrade it to district heating supply temperature. Coefficients of performance (COP) under typical operating conditions range from 2 to 4, depending on the temperature lift between source and supply \cite{david_heat_2017}, \cite{bach_integration_2016}. Capital costs of heat pumps are substantially higher than electrode boilers. The Heat Roadmap Europe \cite{david_heat_2017} survey large-scale installations across European DH systems and document capital expenditure in the range of $300-600 \frac{\euro}{kW_{thermal}}$. The paper by Bach \cite{bach_integration_2016} corroborate the technical performance characteristics for the Greater Copenhagen installations, reporting demand-weighted COP averages of $2.9-3.2$ at distribution level and $2.5-2.6$ at transmission level, alongside full-load-hour profiles of between $3500-4000$ hours per year on the distribution network.

The lower electrical input per unit of thermal output relating to electrode boilers, reduces price exposure. The heat pump's contribution to flexibility is constrained more by the economics and mechanical limits of cycling than by ramp rate; As the temperature lift between source and supply widens, COP deteriorates, and frequent start-stop operations accelerates compressor wear. So while electrode boilers can provide instantaneous up- and down-regulation, large-scale heat pumps have generally not been designed for frequent flexible regulation. \cite{david_heat_2017}. In practice, large-scale heat pumps function as base-load or mid-merit heat sources in the DH merit order, contributing to cost optimization over longer dispatch horizons rather than providing fast-response capacity for services like mFRR activation.\cite{kirkerud_power--heat_2017}, \cite{monie_power--heat_2022}. 

The IEA HPT Annex 57 final report quantifies this tiering through the Sdr. Felding pilot in Denmark, in which a $1.3 MW$ heat pump records a start-up time of approximately seven minutes, and a turn-down time of around four minutes - placing heat-pump participation in the aFRR window rather than the faster, seconds-scale windows accessible to electrode boilers, the Frequency Containment Reserve (FCR) and Fast frequency Reserve (FFR). \cite{v_pedersen_annex_2024}.

Combined heat and power plants (CHP), including waste-to-energy (WtE) incineration facilities occupy a structurally different position within the DH flexibility stack. Rather than being dispatchable against an electricity price signal, they operate in a heat-led mode. Thermal output is the primary production target, driven by waste throughput requirements or heat demand contracts, and electricity is generated as a byproduct at a fixed or minimally adjustable heat-to-power ratio \cite{beiron_combined_2020}. Flexibility provision from CHP and WtE is accordingly indirect, achieved through two mechanisms: adjusting the heat-to-power extraction ratio within the operating range of the turbine cycle, and buffering variable heat production against demand using thermal storage to decouple generation from consumption in time \cite{wang_investigation_2019}. 

Waste incineration supplies roughly half of all Norwegian district heating - the highest share amongst the Nordic countries \cite{sandberg_framework_2018}. In Oslo, Hafslund Celcio's Klemmetsrud WtE plant represents the dominant base-load source in the system merit order, operating continuously with thermal output driven by waste incineration throughput rather than electricity market signals. This heat-led dispatch logic means that WtE capacity constrains the flexibility stack rather than contributing to it directly. Its inflexibility is compensated by the electrode boiler, heat pump capacity and accumulator tank storage. 

\autoref{table2.2:p2h-comparison} summarises the key operational attributes of the three technology classes as they bear on the bedriftsøkonomisk dispatch model.
%

\begin{table}[h!]
\centering
\small
\begin{tabular}{>{\raggedright}p{3.0cm}
                >{\raggedright}p{2.2cm}
                >{\raggedright}p{2.2cm}
                >{\raggedright}p{2.0cm}
                >{\raggedright\arraybackslash}p{3.8cm}}
    \toprule
        \textbf{Technology} &
        \textbf{Capex} &
        \textbf{Opex driver} &
        \textbf{Ramp rate} &
        \textbf{Flexibility role} \\
    \midrule
        Electrode boiler &
        Low &
        Electricity price &
        Near-instant &
        Spot arbitrage; mFRR activation \\[4pt]
        Large-scale heat pump &
        High &
        Electricity price (via COP) &
        Minutes (cycling limits) &
        Base/mid-merit heat source; constrained flexible operation \\[4pt]
        CHP / WtE &
        Very high &
        Fuel / waste throughput &
        Slow (thermal inertia) &
        Indirect via thermal storage buffering \\
    \bottomrule
\end{tabular}
\caption{Operational attributes of the three heat-supply technology
classes relevant to flexibility dispatch in district heating. Electrode
boilers and heat pumps are the power-to-heat (electricity-consuming)
assets \cite{bloes_power--heat_2018}, \cite{david_heat_2017}, \cite{bach_integration_2016}, \cite{monie_power--heat_2022}, \cite{gronli_hvor_2025}, \cite{beiron_combined_2020}.}
\label{table2.2:p2h-comparison}
\end{table}


\subsubsection{Thermal Energy Storage}

The dedicated thermal storage identified in \utsett{AUTOREF section 2.1.1} as the supply-side FTS buffer is what converts the power-flexibility action into a resource that can be shifted across time. It spans a wide envelope of power and duration, from short-duration accumulators suited to day-ahead arbitrage to seasonal store that decouple summer collection from winter demand. Where an asset sits on that envelope fixes which flexibility products it can bid into, following the buffer-parameter logic of \utsett{AUTOREF section 2.1.1}

\utfyll{What is the daily heat delivery of Oslo's district heating compared to the accumulation tank size aka How long can the tank substitute for heat production?}

Thermal accumulator tanks store pressurized hot water at district heating supply temperature inside insulated vessels. They are considered inexpensive amongst energy storage options, as heat water tanks costs between 50 and 100 times less per unit of stored energy compared to electrical storage \cite{hennessy_flexibility_2019}. Their round-trip thermal efficiency is high, on the order of $90-97\%$ under operating conditions; the Danish Energy Agency technology catalogue reports a full-cycle round-trip efficiency of $98\%$ (uncertainty band $96-99\%$) and a standby loss of about $0.2\%$ per day for large-scale steel accumulator tanks, so the operating-condition range used here is consistent with, and conservative relative to, that benchmark \cite{danish_energy_agency_energy_2025}. 

As short-duration, high-power assets, accumulators charge and discharge at rates approaching the connected P2H plant capacity over timescales of hours, making them directly usable for spot-market arbitrage and for buffering the heat produced or displaced during an mFRR activation. The connected electrode boiler can deliver across the full activation window rather than only as long as the heat demand profile permits. Oslo's Haraldrud facility, commissioned in October 2023, is among Norway's largest operational DH accumulators, with a capacity in the order of $300-350MWh$. \avklar{Confirm Haraldrud's exact capacity and its charge/discharge rate. Comparable tanks (Heimdal, trondheim) can frame magnitude}

Accumulator tanks are common in Danish, Finnish, and Swedish district heating but have historically been more scarce in Norway \cite{sandberg_framework_2018}. Outside Oslo, large-scale accumulator capacity in the Norwegian systems is limited, a technical asymmetry treated as a scope boundary in \utsett{AUTOREF section 3.7}

Seasonal thermal energy storage technologies, principally pit thermal energy storage (PTES) and borehole thermal energy storage (BTES), operate on charge-discharge cycles of weeks to months, enabling the temporal decoupling of summer heat collection or waste-heat availability from winter heating demand. The four canonical underground storage typologies - tank, pit, borehole and aquifer - differ substantially in energy density. Tank TES holds roughly $60-80\frac{kWh}{m^3}$, gravel-water PTES roughly $30-50\frac{kWh}{m^3}$, BTES around $15-30\frac{kWh}{m^3}$, and $30-40\frac{kWh}{m^3}$ of stores thermal energy for ATES \cite{dahash_advances_2019}. Round-trip efficiencies are operationally lower than short-duration accumulators and depend strongly on storage utilization cycle; \cite{pan_long-term_2022} report a five-year-average measured storage efficiency of 90.1\% for the Dronninglund PTES at a high utilisation cycle of 2.16, against approximately 66\% for the Marstal PTES operated at a cycle close to 1.0. Power ratings are modest relative to storage volume, limiting participation in balancing markets. Denmark has deployed large-scale PTES extensively, with solar-coupled installations such as Vojens ($200.000 m^3$) and Dronninglund ($60.000 m^3$) \cite{dahash_advances_2019}, \cite{pan_long-term_2022}.

The levelized cost of stored heat at the $100.000 m^3$ scale runs from $\ 130\frac{\euro}{MWh}$ for an accumulation tank to about $160 \frac{\euro}{MWh}$ for the shallow pit, under European cost conditions \cite{dahash_techno-economic_2021}. In Noway, seasonal DH storage remains primarily at the demonstration stage. Kauko document a neighborhood-scale case study at Furuset in Oslo in which district alone reduces peak electricity demand by approximately $28\%$ relative to direct electric heating, while adding borehole seasonal thermal storage reduces peak district heating demand further by around $31\%$, and cuts $CO_2$ emissions by $20\%$ below the direct electric heating baseline \cite{kauko_assessing_2022}. The seasonal-storage configuration is not cost-saving under baseline conditions, however: it carries the lowest operational cost but the highest capital cost, leaving total system cost approximately 3\% above the electric-heating reference. It becomes competitive only when winter electricity prices double or when available grid capacity is reduced by more than $9 MW$ below the site's baseline of $22.6 MW$, a capacity the authors describe as unusually high and present as a special case rather than a representative neighborhood. These figures are therefore indicative at neighborhood scale and do not transfer directly to system scale. The commercial case for seasonal storage is constrained by long capital recovery periods relative to the spot-arbitrage revenues available in the NO1 price environment: unlike accumulators, PTES and BTES assets cannot recover capital cost through mFRR activation or short-horizon arbitrage alone, positioning them primarily as a socio-economic assets in the valuation framework of \utsett{AUTOREF: Section 2.5}. Thorvaldsen, in \cite{thorvaldsen_long-term_2024}, quantify this dynamic for a Norwegian residential setting under the Norwegian monthly-peak grid tariff structure: a strategy that co-optimizes seasonal thermal storage with the monthly-peak capacity tariff achieves a $4.3\%$ total-cost reduction relative to a tariff-only baseline, whereas a storage-only strategy that ignores the tariff signal produces a $69\%$ higher annual cost than the co-optimized case. The mechanism transfers in direction to the DH operator setting: tariff-unaware storage operation systematically under-performs, and the value of storage depends on accessible flexibility products and tariff signals rather than on stored capacity alone.

\illu{Storage typology plot (power in MW on y-axis, duration in time on x-axis) - accumulator tank, BTES, PTES, pipeline storage.}

The value attributed to any thermal storage asset is bounded by the flexibility products available for the stored energy to be sold into, the enabler layer of \utsett{AUTOREF: Section 2.1.1}. Jaanto, in \cite{jaanto_thermal_2023} examine thermal storage deployment under current European electricity market design and find that storage value is limited to the realized price spread across accessible trading windows: in the absence of reserve-market access, accumulators are confined to spot-arbitrage revenue, and seasonal storage captures even less due to low power ratings and poor temporal alignment with daily price signals. This finding applies directly to the Norwegian DH context. Without a flexibility product that compensates the district heating operator for system services beyond mFRR capacity payments, the commercial value of storage investment is anchored on NO1 spot-price variance, \utsett{Price-spread measurement/calculation in Section 2.1.2. - from Nordpool spot price analysis} The societal value of the same storage assets, in the form of avoided grid reinforcement, freed balancing capacity, and system-price suppression, accrues to Elvia, Statnett, and consumers rather than to the operator. The relationship between storage value and market design architecture is therefore a structural feature of the BØ/SØ gap, not a transient market condition.


%--------------------------------------------------------------
%--------------------------------------------------------------
\subsection{Barriers to Flexibility in District Heating}

The corpus on district heating flexibility converges on a recurring set of mechanisms, spanning tariff design, market-product gaps, regulatory frictions, and institutional arrangements, that together separate the operators economic return on flexibility provision from its system-level societal value. The most systematic ordering of these mechanisms from \cite{sneum_barriers_2021}, categorizes them into a taxonomy of forty barriers grouped into nine categories, namely operational signaling, investment, permitting, ownership, district-energy technology conditions, grid access, physical environment, bounded rationality, and acceptance. This thesis collapses those nine categories into the four barrier types that organize the subsections below, mapping district-energy technology conditions, physical environment, and grid access onto technical barriers -- 2.3.1 --
operational signaling and investment onto economic and market barriers -- 2.3.2 --
permitting and ownership onto regulatory and policy barriers --2.3.3--
and bounded rationality and acceptance onto institutional barriers -- 2.3.4--. \cite{sneum_barriers_2021} identifies the absence of a signal-providing scheme as the single barrier sufficient on its own to suppress flexibility, and finds that barriers of national origin dominate the taxonomy, which places the primary responsibility for removing them with national policy-makers rather than with the individual operators. In the resource-and-enabler terms of \utsett{AUTOREF: section 2.1.1}, every barrier in this taxonomy is a missing or adverse \textit{enabler}. The flexibility resource exists physically in the operator's electrode boilers, heat pumps, and accumulator, but markets, reserve products, and network tariffs either fail to reward its activation or actively penalize it. Degefa et al. and Sæle et al. make this split explicit, classifying markets, regulation, and tariffs as enablers separate from the flexibility resource itself \cite{degefa_comprehensive_2021}, \cite{saele_understanding_2023}. The four barrier types below are therefore a catalogue of the enablers that are absent or misaligned in the Norwegian case. \utsett{autoref - section 2.3.5} then sets the Norwegian condition under each barrier type against Denmark and Sweden, the two Nordic systems against which Norwegian district heating is most often measured.

%--------------------------------------------------
\subsubsection{Technical Barriers}

The most visible technical constraint on district heating flexibility provision in Norway is the limited deployment of large-scale thermal storage. The Haraldsrud accumulator, commissioned in October 2023, is Oslo's principal dedicated DH storage asset. Large-scale accumulator capacity elsewhere in the Norwegian system remains limited, confining many operators to within-hour load modulation rather than the multi-hour to day-ahead dispatch optimization that storage enables \cite{hennessy_flexibility_2019}, \cite{sandberg_framework_2018}. Hafslund Celsio's Oslo system, with both a large P2H fleet and dedicated storage, represents a ceiling of Norwegian DH flexibility potential rather than a median, so findings from the Oslo case do not transfer directly onto other Norwegian operators without adjustments in their physical flexibility assets. 

New P2H capability additions face grid connection constraints that delay deployment and cap installed capacity below the economically optimal level. The broad electrification push across sectors has loaded distribution-grid networks and expanded connection queues. In Oslo, applications for roughly $1000MW$ of new electricity load stand against only about $200MW$ of reserved grid capacity, so a DH operator seeking to add electrode boiler capacity may face multi-year connection timelines and capacity restrictions \cite{Thema_fjernvarme_2024}. This caps the P2H capacity an operator can install below economically optimal levels, independent of the downstream market-access question.
\avklar{Verify whether this directly impacts P2H. Verify the "economically optimal levels"}

Network thermal dynamics, including supply temperature frameworks, pipe thermal inertia, and distribution-side control constraints, impose physical limits on the ramp rates achievable through network-level flexibility activation. These are notes as a boundary condition of the economic feasibility in \utsett{AUTOREF: Section 3.7}, but are not subject to detailed analysis in this thesis \cite{vandermeulen_controlling_2018}, \cite{guelpa_demand_2021}.

Across the literature, technical barriers are consistently treated as the most tractable of the four barrier types. Of the forty barriers categorized by \cite{sneum_barriers_2021}, only two concern thermal storage - which the author characterizes as a "no-regrets technology", often seen as the default go-to investment to increase system performance. The bulk of the gap lies in the economic, regulatory, and institutional categories that follow \cite{sneum_barriers_2021}.


%-------------------------------------------------------
\subsubsection{Economic and Market Barriers}

The single most cited economic barrier in the Norwegian DH flexibility corpus is the capacity charge, the monthly peak-load capacity component of Elvia's distribution grid tariff. The capacity charge is assessed on the highest measured power draw within each calendar month, regardless of when that peak occurs or what system conditions prevailed at the time. The regulator itself has acknowledged that this settlement basis can yield a marginal network cost that does not reflect actual network costs, when the monthly peak falls in hours with available capacity \cite{eriksen_proposed_2022}, \cite{Thema_fjernvarme_2024}.  RME corrected this distortion for small low-voltage customers by moving them to a daily-peak settlement. It retained the monthly-peak settlement for the larger customer class, in which a district heating operator the size of Hafslund Celsio sits \cite{eriksen_proposed_2022}. For a district heating operator that uses electric boilers to store thermal energy during low-price hours, any load increase triggers a new monthly peak and incurs a capacity charge proportional to its magnitude. This holds even when the same load increase reduces the aggregate system peak and defers grid reinforcement at the network level. The tariff thereby penalizes precisely the flexibility behavior it should incentivize. For a Norwegian DH plant, electrode-boiler utilization varies between $2\%$ and $17\%$ of annual operating hours depending solely on the grid tariff structure applied, before any market-access barrier is considered \cite{kirkerud_impacts_2016}. The mechanism generalities beyond Norway, though the Norwegian form is unusually severe. For the Danish system, conventional grid tariffs can "obliterate the price signals from the wholesale electricity market \cite{skytte_design_2017}. In that 2014 Danish case, an electric boiler is the cheapest heat source in only about $6\%$ of annual hours, a share that rises to roughly $74\%$ under a flat capacity tariff. The Norwegian capacity charge is a monthly-peak capacity charge rather than the volumetric per-kWh structure analyzed there. The mechanism therefore differs in form and is arguably stronger: the monthly-peak settlement injects intra-month decision uncertainty on top of the constant wedge a volumetric tariff imposes. The direction of the distortion is nonetheless identical. Across the Nordic system, volumetric grid tariffs suppress a substantial volume of cross-sector flexibility investment: a revenue-neutral capacity tariff would unlock on the order of a $20\%$ increase in power-to-heat capacity and a $7\%$ increase in storage \cite{bergaentzle_cross-sector_2022}. Hafslund Celsio is connected at distribution level and faces Elvia's capacity charge on top of the transmission tariff modeled in that study, so the suppression of Oslo-specific flexibility value may be larger still. A comparable conclusion holds at plant level: a capacity-component tariff effectively eliminates electric-boiler operation in Norway and Finland, because a few low-price boiler hours trigger a monthly charge that outweighs the fuel saving \cite{sneum_economic_2018}. Norway's one regulated escape, disconnectable tariff, lowers the capacity charge but requires the operator to accept prolonged disconnection on demand. The discounted tariff is therefore not freely usable for continuous flexibility provision.

This does not make the capacity charge a barrier simply to be abolished. Capacity-charge designs are not inherently barrier-forming. Removing the capacity component while retaining a dynamic energy signal lets assets with high thermal mass pre-heat ahead of high-price periods, which simply shifts peaks rather than reducing them \cite{wilczynski_assessment_2023}. Some capacity signal is therefore needed for grid-peak suppression. The barrier lies in the specific peak-based monthly structure of the Norwegian capacity charge, not in capacity pricing. The design flaw is in which peak the tariff targets. On a real Oslo case, a capacity-based tariff alone incentivizes individual peak reduction without coordinating the coincident peak that actually drives grid capacity needs \cite{askeland_activating_2021}. Closing the gap between individual and coincident peak behavior therefore requires a complementary local capacity-trading mechanism that the Norwegian system currently lacks. Oslo demand-response data corroborate this: the individual-peak, ex-post settlement of the capacity charge performs worst for substation-level peak reduction, even as it performs best at the individual-customer level \cite{hofmann_grid_2025}. The tariff nonetheless materially governs storage-coupled dispatch. Co-optimizing a Norwegian building's seasonal-storage operation against the monthly measured-peak tariff, by absorbing cheap summer electricity to lower the costlier winter peak charge, reduces annual operating cost by $4.3\%$ and the monthly-peak tariff bill by $23.9\%$ relative to operation without storage \cite{thorvaldsen_long-term_2024}. The residential scale of these cases does not transfer numerically to the district heating operator. The mechanism, however, applies directly to Hafslund Celsio's electrode-boiler and accumulator dispatch under the capacity charge. Finally, the capacity charge does not bind all assets equally. A monthly-maximum load-demand tariff of the Oslo type is the worst tariff for electric-boiler deployment but the best for heat pumps: base load heat-pump operation spreads the capacity charge across many hours, whereas intermittent boiler dispatch concentrates it into a few. Tariff reform therefore releases more flexibility from electrode boilers than from heat pumps, the latter remaining base load regardless of tariff design unless a competing low-cost heat source emerges \cite{sandberg_energy_2019}, \cite{johannsen_incentivising_2021}.

Until early 2025 the mFRR minimum bid threshold functioned as a structural exclusion mechanism, converting the cost asymmetry above into a permanent revenue gap. At $10 MW$ nationally and $5 MW$ in NO1 and NO3, the floor placed reserve-market revenue out of reach for all but the largest aggregated capacity. As established in \utsett{AUTOREF section 2.1.2}, the Nordic harmonization of the mFRR bid floor to $1 MW$ on 4 March 2025 removes the bid-size barrier as such. A single Hafslund Celsio electrode-boiler site now qualifies on scale alone. What remains is a revenue gap of a different origin. For the small plants that compose the sub-threshold fraction of the national fleet, prequalification cost, baseline and metering complexity, and the absence of any aggregation framework to pool capacity still foreclose participation \cite{monie_power--heat_2022}, \cite{gronli_hvor_2025}. The reserve products also differ in how well district heating assets match them, so participation turns on asset capability rather than scale alone. Fast-ramping electrode boilers can serve the frequency-containment products, whereas the slower heat-pump units are confined to the automatic restoration (aFRR) and mFRR windows that respond on a timescale of minutes \cite{javanshir_operation_2022}, \cite{boldrini_role_2022}. THEMA quantifies the adjacent system-side magnitude: district heating conversion in Oslo would free roughly $250-430 MW$ of grid capacity, against Hafslund Celsio's installed electric-boiler fleet of about $283 MW$ \cite{Thema_fjernvarme_2024}.

The reserve product exclusion is compounded by the absence of any market instrument through which district heating operators can be compensated for system services that fall outside the mFRR activation window. This is a missing-market problem \cite{mitridati_heat_2020}. Because sequential heat-then-electricity clearing is short sighted, the district heating system provides storage-equivalent - "Virtual electricity storage" - services to the power system through synchronized P2H dispatch. Without a bilateral flexibility contract or a purpose-designed capacity product, these services go effectively uncompensated. No Norwegian market mechanism currently offers a capacity payment for DH-side load flexibility beyond the mFRR activation payment itself. The absence of such a product is confirmed across multiple Nordic studies \cite{sneum_economic_2018}, \cite{jaanto_thermal_2023}, \cite{selvakkumaran_how_2021}. The THEMA report attributes the gap to specific regulatory-design choices rather than to a textbook market failure: district heating receives no payment for the grid-service value it provides, the regulated price cap is computed monthly rather than hourly so hourly flexibility revenue is blunted, and operators in effect self-finance the \utsett{How to translate strømstøtte} scheme for their customers \cite{Thema_fjernvarme_2024}. THEMA's headline estimate of NOK 1-1.6 billion per year as the societal value of converting heat supply to district heating, that is, fuel substitution and peak reduction. This is a broader quantity than the flexibility-specific value this thesis isolates. The flexibility subset is correspondingly smaller, and it is the quantity \utsett{AUTOREF Section 4.3} attributes across the barrier mechanisms.
\avklar{VERIFY: the NOK 1-1.6 billion/year figure is THEMA's CONVERSION value (fuel substitution + peak reduction), NOT the flexibility value. The flexibility component is far smaller (the FCR-N illustrative NOK 330,000/MW/yr line). keep §4.3 anchored on the flexibility subset. Thema = Sponsored report; corroborate with other sources.}

The NO1 spot price structure sets a ceiling on the arbitrage revenue available to Oslo-based operators, independently of market-access barriers. As established in \utsett{autoref sec 2.1.2}, NO1 exhibits structurally lower price-duration variance than NO2 and NO5, which compresses the price spread that drives P2H dispatch value. Capacity-based tariffs of the Norwegian type reduce effective electrode-boiler utilization to low levels \cite{sandberg_energy_2019}. Combined with the compressed NO1 price spread established in \utsett{autoref sec 2.1.2}, this means the flexibility investment cannot recover on commercial returns alone in the Norwegian context. That reinforces the case for complementary reserve-market access rather than sole reliance on spot-arbitrage revenue.

Beyond the regulated price ceiling itself, the price models through which Norwegian operators bill heat do not transmit flexibility value to customers. Conventional volume-based district heating price models are inadequate for a flexible supply-and-demand system, and price-model design directly governs an operator's ability to capture flexibility value \cite{lygnerud_capturing_2024}. The time-differentiated models that have supported flexibility investment elsewhere in the Nordic market are foreclosed in Norway, both by this price-model convention and by the alternative-cost ceiling in energy law §5-5 examined in \utsett{autoref Section 2.3.3}. An operator that reduces system peak load and contributes to lower electricity prices therefore cannot recover that value through the heat price, because the ceiling adjusts downward with the electricity price it helped reduce.


%-------------------------------------------------------
\subsubsection{Regulatory and Policy Barriers}

District heating prices in Norway are governed under §5-5 of the Energy Act. The wording originates in the 1986 District heating act and was carried forward into the 1991 Energy act, so there is no separate laws governing district heating \cite{nve_nves_2024}. The provision sets a ceiling on DH tariffs at the cost of the nearest competitive heating alternative for the customer. This cap-style mechanism is designed to protect customers from monopoly pricing, but it structurally suppresses any flexibility-value signal that would otherwise flow through the heat price. The THEMA report identifies this regulatory framework as a primary source of misaligned incentives \cite{Thema_fjernvarme_2024}. An operator who invests in flexibility assets that reduce system peak load and enable lower electricity prices at the system level cannot recover that investment through higher heat prices, because the price ceiling adjusts downward with the electricity price the operator helped reduce. DH price regulation anchored below long-run supply cost produces structural underinvestment across European DH systems \cite{odgaarda_review_2020}. A concrete Danish quantification of this dynamic under tariff-reform conditions is provided in \cite{djorup_district_2020}.

The electricity tax applied to P2H consumption constitutes a cost component that reduces the operating margin of electric boilers and heat pumps. The Norwegian industrial rate is 0.60 øre/kWh in both 2025 and 2026, broadly comparable to the Danish industrial rate of approximately 0.80 øre/kWh \utfyll{source: A3: Årsrundskriv for avgift på elektrisk kraft 2026} \cite{Thema_fjernvarme_2024}. The narrow gap between the two countries and the already-low absolute level mean that further tax reduction has limited scope as a flexibility policy lever. Relative to the capacity charge and the mFRR minimum bid, the electricity tax is a second-order barrier, treated in the commercial model as a fixed operating-cost component rather than a primary gap mechanism. A historical precedent underscores the point. The Danish electricity-to-heat tax of 67.4 €/MWh had to be zeroed out before Copenhagen district heating electric boilers became viable for wind-power integration \cite{meibom_value_2007}. An input-side cost wedge on power-to-heat can therefore block flexibility outright when it is large. In the Norwegian case the binding wedge is no longer the tax, which is small, but the capacity charge capacity charge analyses in AUTOREF: Section 2.3.2.

A Norway-specific regulatory distortion is the $CO_2$ levy applied to waste incineration. For 2025 the levy is set at NOK 908 per tonne of $CO_2$ for non-quota-liable emissions and NOK 182 per tonne for quota-liable (EU ETS) emissions. For 2026 the non-quota rate falls to NOK 848 per tonne, while ETS-covered plants are exempted \cite{noauthor_stortingsvedtak_2024}. \avklar{830 NOK} This levy has no direct counterpart in Denmark or Sweden. It directly affects the operating economics of WtE-anchored DH systems such as Hafslund Celsio's Klemetsrud plant, by raising the marginal cost of the dominant heat source in the merit order. The applicable rate depends on Klemetsrud's quota status, and the planned carbon capture at the plant (operational from about 2029) would exempt the captured share entirely. A rising WtE cost shifts the dispatch balance toward electric boilers and heat pumps, potentially increasing the economic attractiveness of P2H flexibility. It simultaneously compresses operator margins and may reduce the financial headroom for the capital investment that flexibility requires. The net effect on the commercial/socio-economic gap through this channel is directionally ambiguous, and is flagged as an open analytical question in AUTOREF: Section 5.5.

The extension of \utsett{Oversetting for "Norgespris"} to district heating customers represents a compounding policy barrier. The Norgespris program, originally designed to cap residential electricity bills, was extended to DH customers with effect from 1 October 2025 at a fixed 50 øre/kWh including VAT, and is state-funded from that date \cite{nve_norgespris_2026}. Prior to the extension the operator self-financed the \utsett{Oversetting for strømstøtte} equivalent for its customers. Hafslund Celsio's audited accounts record cumulative district-heating losses on the order of NOK 1 billion over the preceding three years \cite{hafslund_arsrapport_nodate}. The regime change therefore shifts the burden from the operator's profit \& loss to the state budget. The price-ceiling mechanism in Energy act § 5-5 \utsett{add Energy act source }continues to apply, now indexed to a politically set retail electricity price rather than to a market price. The THEMA NOK 1-1.6 billion conversion-value figure is based on pre-Norgespris conditions. The post-October 2025 regulatory configuration removes one operator-side cost channel while preserving the underlying flexibility-value-capture gap, since the § 5-5 ceiling still suppresses any flexibility-driven margin uplift the operator might otherwise realize.

%--------------------------------------------------------
\subsubsection{Institutional Barriers}

The most consequential institutional gap is the absence of a formal aggregation framework through which multiple sub-threshold DH operators could pool flexibility capacity to clear mFRR minimum bids collectively \cite{sneum_barriers_2021}, \cite{ma_literature_2020}. The Swedish SthlmFlex pilot, run jointly by Svenska kraftnät, Vattenfall Eldistribution, and Ellevio from winter 2020/2021, demonstrated the technical and institutional architecture of multi-operator local flexibility coordination. It included a reduced mFRR minimum bid of 1 MW for sthlmflex-linked bids, against the standard Swedish 5 MW threshold \cite{nyman_how_2026}. The pilot has since concluded with very low market liquidity. Ellevio's own assessment is that short-term LFMs are not a viable congestion-management instrument for Stockholm in the near term \cite{scrocca_local_2026}. The Swedish experience therefore demonstrates both the institutional template and its commercial fragility. Even with reduced bid thresholds and explicit TSO-DSO coordination, sellers found revenue insufficient to justify participation; one aggregator estimated LFM income at approximately $10\%$ of ancillary-market income for equivalent effort \cite{nyman_how_2026}. The Danish EcoGrid 2.0 field trial ( 800 residential customers on Bornholm, 36 real activations) provides complementary peer-reviewed evidence that the LFM mechanism is operationally feasible at scale \cite{heinrich_ecogrid_2020}. The qualification is that flexibility services act as an insurance product against rare-but-consequential overloading events, rather than as a continuous revenue stream. Twenty-five interviews across the parallel CoordiNet Uppland and Skåne LFM pilots corroborate this barrier inventory, finding that costs exceeding benefits, lack of automation, and market complexity outweighed monetary incentives as the dominant participation barriers \cite{palm_drivers_2023}. A Norwegian DSO interview study reaches the same verdict from the buyer side. Distribution companies rate local flexibility markets below the prominence they receive in the academic literature, and report that qualitative concerns, namely availability, predictability, controllability, and operational risk, weigh more heavily in their adoption decisions than cost-benefit calculations \cite{saele_understanding_2023}. In Norway, no equivalent framework exists, so the mFRR threshold exclusion documented in Section 2.3.2 is institutionally locked in rather than a temporary commercial gap. Across a systematic review of 38 district-energy business-model studies, the absence of aggregating intermediary institutions is a recurring structural barrier to flexibility-market participation in DH systems globally \cite{selvakkumaran_how_2021}. Norway's position exemplifies this pattern.

A second institutional gap is the structural disconnect between the distribution system operator (DSO) and the transmission system operator (TSO) in the attribution of flexibility value. Flexibility provided by Hafslund Celsio's electric boilers reduces peak load on Elvia's distribution network, deferring grid reinforcement investment and accruing value to Elvia and ultimately to grid users through lower tariffs. Simultaneously, mFRR activation of the same assets provides balancing services to Statnett at the transmission level. Neither value flow reaches the district heating operator through any existing mechanism. No Norwegian local flexibility market analogous to SthlmFlex bridges the distribution-level gap, and the TSO-level mFRR market excludes most operators through the minimum bid threshold \cite{brolin_design_2020}. The result is that district heating flexibility assets operate in a dual-benefit landscape where neither institutional beneficiary compensates the provider.

A third institutional condition specific to the Oslo case is Hafslund Celsio's position within the Hafslund group, which holds assets across hydropower generation and district heating and a regional grid stake (Elvia), under City of Oslo ownership \cite{hafslund_interim_nodate}. In principle, vertical integration of this type could align incentives for flexibility by internalizing the grid value that DH flexibility creates within the same corporate perimeter. In practice, capital-allocation decisions within a diversified energy group respond to group-level return requirements and portfolio priorities that may not reflect the marginal socio-economic value of district heating flexibility investment \cite{fernqvist_district_2023}. This is presented as a structural condition of the Oslo case study, not as an evaluative judgment on Hafslund's decision-making. Its practical significance for the flexibility gap is an open question that the operator interview in Section 3.5 is positioned to explore.


%--------------------------------------------------------------------------
\subsubsection{Cross-Country Comparison: Norway vs Denmark and Sweden}

The four preceding subsections have characterized each barrier type within Norway's regulatory and market context, deliberately keeping Danish and Swedish material to the periphery to preserve the taxonomy structure. This section draws those threads together in a three-country synthesis. Sweden is treated as an equal comparator to Denmark, correcting the under-representation flagged in the Consensus synthesis and documented in the Swedish source literature \cite{fernqvist_district_2023}, \cite{atabaki_optimising_2026}.

\utsett{AUTOREF Table below (crosscountry Nordic)} compares Norway, Denmark, and Sweden across the seven dimensions on which regulatory and market conditions most directly shape the commercial DH value of district heating flexibility.

\todo{Metric locked to national share of building heat demand, applied uniformly: NO  10-12\% (Sandberg 2018), DK  50-65\%, SE  50\% (Werner 2017). City use the city's share of its own heat demand (Oslo  20\%, OED 2022 via NOU 2023:3); dobbeltsjekk DK/SE city figures with the same basis (national share of building heat demand).}

\utsett{All sources for this: }
\begin{table}[h!]
\centering
\caption{Framework conditions for district heating flexibility:
Norway, Denmark, and Sweden compared. Sources: \cite{sandberg_framework_2018}, \cite{sneum_economic_2018}, \cite{kirkerud_impacts_2016}, \cite{johansen_something_2022}, \cite{fernqvist_district_2023}, \cite{atabaki_optimising_2026}, \cite{levihn_chp_2017}, \cite{sandberg_energy_2019}, \cite{skytte_flexible_2019}, \cite{Thema_fjernvarme_2024}}
\small
\begin{tabular}{>{\raggedright}p{3.2cm}
                >{\raggedright}p{3.5cm}
                >{\raggedright}p{3.5cm}
                >{\raggedright\arraybackslash}p{3.5cm}}
\toprule
\textbf{Dimension} &
\textbf{Norway} &
\textbf{Denmark} &
\textbf{Sweden} \\
\midrule
Grid tariff structure &
monthly peak-based capacity charge; penalizes flexible load regardless of system benefit &
Time-of-use design; progressive reform trajectory supporting P2H dispatch &
Capacity-based components; SthlmFlex introduces local tariff flexibility at distribution level \\[4pt]
mFRR minimum bid &
1 MW Nordic-harmonized since March 2025 (was 10 MW national / 5 MW NO1 \& NO3); bid size no longer the binding exclusion, but no aggregation framework &
1 MW (Nordic-harmonized); cooperative / BRP structures ease aggregation &
1 MW (Nordic-harmonized); SthlmFlex pilot trialed local aggregation (since concluded) \\[4pt]
Spot price variance (P2H arbitrage headroom) &
Low (NO1); structurally compressed arbitrage margin &
High (DK1/DK2); large P2H arbitrage window driven by wind integration &
Moderate; improving with wind capacity build-out \\[4pt]
DH penetration, national / capital city &
$\sim$10--12\% / $\sim$20\% Oslo$^{a}$ &
$\sim$50--65\% / $\sim$98\% Copenhagen$^{a}$ &
$\sim$50\% / $\sim$50\% Stockholm$^{a}$ \\[4pt]
Aggregation framework &
None; no formal mechanism to pool sub-threshold operators &
Cooperative ownership \& balancing responsible party structures enable implicit aggregation &
SthlmFlex local flexibility market pilot; aggregation emerging \\[4pt]
Electricity tax on P2H (industrial rate) &
$\sim$0.60 øre/kWh &
$\sim$0.80 øre/kWh (broadly comparable; limited policy leverage) &
Comparable; DH-specific exemptions in some contexts \\[4pt]
Sector-coupling governance &
No named policy instrument; ministerial signals of openness (2025) but no formal mechanism \cite{hovland_varsler_2025} &
DH-electricity coupling a named national energy policy priority; multiple instruments in place &
Emerging via national energy planning and local market pilots \\
\bottomrule
\multicolumn{4}{l}{\footnotesize $^{a}$ Figures are indicative;
see verification note above. Definition: national or city-level
share of building heat demand.} \\
\end{tabular}
\end{table}

The comparison reveals a consistent structural pattern. Norway has not updated its regulatory and market architecture to reflect the changed system conditions in which district heating flexibility has become system-valuable. Denmark and Sweden have, through different institutional pathways, partially closed the gap between the social and private value of DH flexibility \cite{sandberg_framework_2018}, \cite{sneum_economic_2018}. The Norwegian distance from the comparators on any individual dimension is moderate. The grid tariff is more restrictive than Danish tariff structures but not qualitatively different from Swedish capacity charges, and the mFRR bid floor is now identical across the Nordic markets. The distinctive feature of the Norwegian case is that the remaining dimensions stack simultaneously in the direction that suppresses commercial flexibility value, producing a compounding barrier geometry that neither comparator replicates. This compounding interaction is the defining structural characteristic of the Norwegian DH sector's anomalous position within the Nordic group \cite{sandberg_framework_2018}. The barrier taxonomy documented in \utsett{Sections: 2.3.1, 2.3.2, 2.3.3, 2.3.4} operationalizes each layer of that interaction separately, so that \utsett{Autoref section 4.3} can attribute an economic value contribution for each.



%-------------------------------------------------------------------
%-------------------------------------------------------------------
%-------------------------------------------------------------------
\subsection{The Norwegian Context}

\utsett{Subsection 2.4 opening statement}


\subsubsection{Energy System Characteristics}

Norway's electricity system is dominated by hydropower to a degree unmatched in any other European country, with hydropower accounting for approximately $90\%$ of annual electricity generation and the remainder drawn primarily from wind, supplemented by gas turbines and a small thermal reserve \cite{NVE_langsiktig_2023}, \cite{iea_-_international_energy_agency_norway_2022}. This structural condition creates a marginal generation unit that behaves differently from the thermal or wind-based units that set the price in most European systems. The marginal cost of a Norwegian hydropower reservoir is not a fuel price but an opportunity cost: the shadow value of stored water, determined by expected future scarcity, cross-border export prices, and the scarcity of balancing reserves rather than any single commodity input. As established in \utsett{AUTOREF section 2.1.3}, district heating electric boilers that respond to low spot prices in this system do not primarily displace fuel combustion. They conserve reservoir capacity for higher-value export windows or reserve provision. The socio-economic value pathway, formalized as one of the two valuation perimeters in \utsett{autoref sec 2.5}, runs through hydro reservoir economics and cross-border market design, a characteristic with no direct parallel in fossil-heavy or wind-dominated European systems. This pathway has been quantified directly \cite{askeland_balancing_2019}. In a reference year, roughly 149 hours see Norwegian hydropower generation capacity, rather than the inter-connector, bind the export limit. A shift of approximately $25\%$ of individual electric heating to district heating removes that binding constraint in every hour, freeing reservoir capacity for export and balancing. The same shift carries a trade-off, reducing Norway's capacity to absorb surplus European renewable generation, so the net system value is conditional rather than unambiguous. Independently of the operator-funded literature, NVE's long-term market analysis identifies district heating electric-boiler coupling as a long-duration flexibility resource, more effective than batteries for multi-hour and multi-day surplus \cite{NVE_langsiktig_2023}. The same analysis projects the Norwegian capacity balance tightening toward a deficit in the most stressed hours by 2040.

\illu{LAGE FIGUR 2.4.1: two panels; physical attributes of the Norwegian DH flexibility system. Panel 1: simplified flow diagram connecting the power system (hydro/wind generation, Statnett transmission, Elvia distribution) to the DH operator's P2H assets (electrode boiler, heat pump), the accumulator, network thermal mass, and building heat demand. Label actor nodes (Statnett, Elvia, Hafslund Celsio). Panel 2: price-duration curve for NO1 alongside DK1 or SE3 (verify), showing the compressed low-price tail.}

The Norwegian power grid is divided into five bidding zones, NO1 through NO5, introduced in \utsett{Autoref sec 2.1.2}, within which a single spot price clears each hour via Nord Pool. Transmission constraints between zones produce systematic price differentials during periods of high cross-border flow or intra-Norwegian congestion. Norwegian spot prices exhibit comparatively low price-duration variance, which the Flex4RES assessment attributes to the system's dispatchable hydropower dominance, the property that makes Norway the most price-stable area in the Nordic region \cite{skytte_flexible_2019}. Oslo and the surrounding eastern grid fall within NO1, a zone with negligible local wind or other variable-renewable generation, so its hourly price profile is shaped by hydropower dispatch, interconnector schedules, and demand variation rather than by intermittent renewable production. For district heating operators based in Oslo, the direct consequence is a compressed arbitrage window: the spread between off-peak and on-peak spot prices is narrow, which reduces the commercial return on power-to-heat flexibility before any market-access barrier is applied. This compression is not a contingent feature of current market conditions but a structural property of the zone's generation mix, reinforcing the gap between commercial and socio-economic value on the revenue side. The same compression measured against Danish and Swedish reference zones is set out in the cross-country comparison in \utsett{autoref sec 2.3.5}.

Norwegian heating demand is structurally concentrated in winter and skewed toward direct electric resistance heating, a legacy of the historically low electricity prices that made direct electric heating economically rational for households and commercial buildings throughout most of the twentieth century \cite{NVE_langsiktig_2023}. Norwegian electricity demand is correspondingly highly price-inelastic at the short horizon. For southern Norway (the NO1, NO2 and NO5 zones), the short-run price elasticity is estimated at approximately $-0.09$ over 2013--2023, a period that includes a more-than-fourfold increase in average spot prices, and it becomes statistically indistinguishable from zero in the coldest, highest-price periods \cite{idso_short-term_2024}. Temperature explains the bulk of demand variation. The price signal that drives operator-side flexibility decisions does not pass through to consumer demand in any short-run manner relevant to the commercial/socio-economic gap framing. This direct electric heating stock creates a sharp winter peak in electricity demand that intensifies both congestion on the NO1 grid and the scarcity value of generation capacity during peak periods. District heating networks can in principle reduce this peak by converting electricity to stored thermal mass during off-peak hours and supplying building heat from storage during morning and evening demand peaks, a function that corresponds to the virtual electricity storage concept introduced in \utsett{ref sec 2.1.3}. The system-level scale of the available resource has been quantified in a stochastic TIMES-Norway model run to 2050: residential space heating accounts for $83-90\%$ of all cost-efficient demand-response operation, dwarfing both industrial and commercial sources \cite{ahang_analysis_2025}. The same model shows demand response smoothing Norwegian electricity prices and reducing regulated-hydropower income by $0-4\%$ across the five price zones, which locates a meaningful share of the socio-economic benefit on actors other than the operator that provides the flexibility. That framing is residential-side and this thesis values flexibility from the operator-supply side, but the dispatch pattern is the same: thermal flexibility shifts winter electric heating load away from peak. In the BALMOREL Nordic context, power-to-heat in flexible district heating and increased transmission capacity "seem to have a larger economic potential than household and tertiary demand response \cite{kirkerud_role_2021}. The extent to which this peak-shaving function is commercially accessible to district heating operators depends on the reserve products and tariff structures available to them, a question addressed directly in \utsett{autoref sec 2.3 and 3.3}.

The Norwegian district heating sector is small by Nordic standards but geographically concentrated. Aggregate national net production was approximately $7568 GWh$ in 2024, of which about $6810 GWh$ was delivered to consumers after distribution losses \cite{boeng_mindre_2025}. By energy source, 2024 net production was roughly $38\%$ waste incineration, $29\%$ solid biofuel, $12\%$ electric boilers and $11\%$ heat pumps, giving Norway the highest power-to-heat share in the Nordic district heating group \cite{boeng_mindre_2025}, \cite{sandberg_framework_2018}.


\todo[inline]{Cross-check: confirm the  10-12\% national DH penetration (Sandberg 2018, used in §2.2.1 and the §2.3.5 table) reconciles with the 2024 SSB volumes above (6810 GWh delivered / 7568 GWh net production) against total national building heat demand; pin the denominator from a consistent SSB basis so the percentage and the absolute volumes agree.}

\illu{FIGURE (§2.4.1): an SSB-data fuel-mix / fjernvarmebalanse graphic from SSB 2024 StatBank tables 09469 (net production by source) and 04727 (gross→net→loss→delivered) would visually anchor the high P2H share — a stacked bar or Sankey of the 2024 mix. Produce externally in Python/Matplotlib per the project figure rule. Drop if the §2.4.1 schematic figure already covers the need.}

Hafslund Celsio's Oslo system delivers approximately 2{,}000 GWh annually, equivalent to roughly 20\% of Oslo's heating demand and 25\% of its capacity demand \cite{NOU 2023:3 Mer av alt - raskere - Energikommisjonen - 2023}.

\avklar{VERIFY: NOU 2023:3 underlying source for the 20\% / 25\% Oslo figures is OED (2022) Har vi fleksibilitet nok — when this framing recurs in §1 / §4 / §5, cite both together.} 

Trondheim and Bergen operate secondary systems under Statkraft Varme and BKK Varme respectively, with the remainder distributed across smaller municipal and cooperative networks. The Oslo system is substantially more technically capable and instrumented than any of the secondary systems, owing to its scale, the diversity of its generation stack (waste-to-energy base-load, electrode boilers, and large-scale heat pumps), and the Haraldrud accumulator commissioned in October 2023.

\avklar{VERIFY: Hafslund Celsio total installed P2H capacity. Company communications cite a headline of 283 MW, but a per-site sum of documented electrode boiler and heat pump installations yields a more conservative estimate of approximately 131-137 MW. Resolve the discrepancy against primary commissioning documentation before finalizing. Use the conservative per-site sum in the model (Section 3.3).}

\avklar{WRITE: confirm Haraldrud's exact capacity (the 300 and 350 MWh figures both circulate) and its peak charge/discharge rate (MW) against Hafslund Celsio commissioning documentation or via interview. Both are required as Section 3.3.2 model inputs.}

The Haraldrud facility is among Norway's largest operational district heating accumulators, set against the other Norwegian thermal stores cataloged in \utsett{autoref sec 2.2.3}. Its significance for this thesis extends beyond its function as a model input. Hafslund Celsio invested in short-duration thermal storage capacity despite the compressed spot-price variance of NO1 and the restricted reserve-product access documented in \utsett{sec 2.3}. This indicates that the operator views the commercial flexibility case as viable on projected rather than only observed market conditions. Large-scale accumulator capacity is unevenly distributed across the Norwegian systems, an asymmetry that limits the quantitative model in \utsett{sec 3.3.} to the Oslo system and is treated as a scope boundary in \utsett{ref Section 3.7}.


%-----------------------------------------------------------
\subsubsection{Policy Evolution 2015-2025}

The policy environment governing Norwegian district heating flexibility has shifted materially over the decade from 2015 to 2025. It accumulated institutional changes that individually addressed specific inefficiencies without resolving the market-design mismatch that is the thesis's central object of study. Tracing this evolution situates the quantitative modeling in \utsett{autoref: Chapter 3} against a specific regulatory timeline, and frames the barrier analysis in \utsett{autoref: Section 2.3} as a snapshot of a system in incremental transition rather than a static condition.

Grid tariff reform entered the regulatory agenda in 2015, when NVE opened its first consultation on replacing the then-current distribution-network tariff with a capacity-based element intended to reduce peak congestion and improve price signals for flexible consumers \cite{andersen_horing_2015}. These consultations ultimately produced the capacity charge structure, a monthly peak capacity charge levied on the average of the three highest hourly power draw within a calendar month. As documented in \utsett{ref: Section 2.3.2}, this regime functions as a commercial barrier to the specific flexibility behavior it was ostensibly designed to encourage. A district heating operator that raises its electrical load during an mFRR activation or a spot-arbitrage event risks setting a new monthly peak that persists for the full billing period, which makes the net financial case for activation ambiguous even when the spot-price signal is favorable \cite{kirkerud_impacts_2016}.

A Nordic-incentives analysis covering Norwegian district heating was published in 2018 \cite{sneum_economic_2018}. The Flex4RES consortium's summary report, a Nordic cross-country assessment of demand-side flexibility barriers that also covered Norwegian district heating, followed in 2019 \cite{skytte_flexible_2019}. The barrier taxonomy later adopted as the scaffolding for \utsett{ref: Section 2.3} is the 2021 taxonomy study \cite{sneum_barriers_2021}. These documents pre-date the Haraldrud commissioning and the 2025 Norgespris extension, establishing a reference baseline against which subsequent regulatory developments can be compared.

The institutional capacity to reform these conditions consolidated over the same period. The Reguleringsmyndigheten for energi (RME) became a formally independent regulator on 1 November 2019 under the European Union's Third Energy Market Package \cite{stortinget_vedtak_2019}, separating market regulation from NVE's broader resource-administration mandate.

The 2022 International Energy Agency review of Norwegian energy policy flagged district heating flexibility barriers among several under-utilized demand-side resources and recommended clearer market access rules and aggregation provisions \cite{iea_-_international_energy_agency_norway_2022}.

The year 2023 brought parallel advances in the evidence base and in the physical system. The SINTEF ZEN research program's "Report 47" modeled the Norwegian building stock to 2050, quantifying how energy efficiency combined with district heating and heat-pump deployment could lower national peak electricity demand and reduce system energy costs \cite{SINTEF_energy_2023}. Vista Analyse published a cost-benefit analysis of energy-flexible heating systems at the building level, assessing the private and socio-economic consequences of possible changes to the TEK17 building-code flexibility requirements \cite{Vista-analyse_energifleksible_2023}. Its valuation of avoided generation and grid investment informs the socio-economic framing in \utsett{autoref: Section 2.5}. In October 2023, the Haraldrud accumulator was commissioned, materially changing the technical basis of Oslo system flexibility.

In 2024, Hafslund Celsio commissioned a market analysis from THEMA Consulting \cite{Thema_fjernvarme_2024}, the first operator-funded external assessment of the revenue gap from the commercial perspective. The THEMA report documents the mFRR capacity revenue shortfall under current market conditions and provides preliminary estimates of the value that an accessible flexibility product could unlock; these estimates serve as reference comparators in \utsett{Ref: sec 4}.

The year 2025 brought two regulatory developments with direct bearing on the commercial case for district heating flexibility. With effect from 1 October 2025, the Norgespris electricity price support scheme was extended to district heating customers and placed on state funding from that date. The heat price passed through to end consumers is now indexed to a politically set retail electricity price rather than to a market price, and Hafslund Celsio's pre-extension self-financing of strømstøtte (which had produced approximately NOK 1 bn in cumulative losses over three years) is replaced by state reimbursement \cite{nve_norgespris_2026}, \cite{hafslund_arsrapport_nodate}.

For 2025 the $CO_2$ levy on waste incineration was set at NOK 908 per tonne for non-quota-liable emissions, and NOK 182 per tonne for quota-liable emissions. The non-quota rate falls to NOK 848 per tonne in 2026, with ETS-covered plants exempted. This reshapes the cost basis of the waste-to-energy base-load that anchors the Oslo merit order \cite{stortinget_avfallsforbrenning_2025}.

The energy minister signaled awareness of the flexibility utilization gap in public statements in late 2025 without committing to a specific new market product \cite{hovland_varsler_2025}. A 2025 consultancy note by Grønli from AFRY offers directional estimates of the mFRR revenue potential available to Norwegian district heating capacity under current and reformed threshold conditions \cite{gronli_hvor_2025}.

A 2026 sector-coupling optimization for the Swedish district heating sector is the most recent structurally analogous study in the Nordic literature \cite{atabaki_optimising_2026}. Its Norwegian parallel at the firm level, with paired commercial and socio-economic valuation, is the gap this thesis addresses


%------------------------------------------------------------
\subsubsection{Why Norway Is Its Own Case}

The preceding subsections establish two structural claims. \utsett{ref Section 2.4.1} documents that the Norwegian power system routes the socio-economic value of power-to-heat flexibility through hydro reservoir economics and cross-border export pricing, a value pathway absent from the European literature that developed primarily around fossil displacement and VRE curtailment avoidance. \utsett{Section 2.4.2} documents that a decade of incremental regulatory activity has modified specific elements of the market environment without resolving the core mismatch between where flexibility value is created and where its cost is borne. This subsection addresses whether these two features, in combination, make Norwegian district heating flexibility analytically distinct from any Nordic comparator, or whether it constitutes a scaled-down version of the Danish or Swedish case that can be modeled by adaptation.

The hydro value pathway is necessary but not sufficient to explain the structural character of the Norwegian case. If it were the only distinguishing feature, the analytical gap could in principle be addressed by adapting a Danish or Swedish sector-coupling optimization and substituting the NO1 price signal for DK1 or SE3. What prevents this reduction is the co-occurrence of four conditions that individually appear elsewhere in the Nordic group but do not stack simultaneously in Denmark or Sweden \utsett{ref Section 2.3.5}: a substitution-price cap on heat revenues under the Energy Act §5-5, under which the heat price may not exceed the cost of alternative electric heating and is therefore now indexed to a politically set retail electricity price since the state-funded \utsett{oversetting for Norgespris} extension took effect for district heating customers on 1 October 2025; a national district heating penetration of roughly $10-12\%$ that structurally limits the scale at which any single operator's flexibility affects system outcomes; an capacity charge tariff structure that penalizes the load-raising behavior that mFRR activation requires; and a $CO_2$ levy trajectory that raises the operating cost of the waste-to-energy base-load anchoring the Oslo merit order. No single Nordic comparator faces all four simultaneously, and the compounding interactions between them, in particular the joint effect of the price cap and the tariff structure on the operator's net margin from a flexibility event, are not captured by any of the single-condition barrier studies reviewed in \utsett{ref Section 2.3}.

The consequence for the valuation framework is that the gap between commercial and socio-economic value of district heating flexibility in Norway is not merely a magnitude difference relative to Danish or Swedish equivalents: it is a structural outcome produced by a specific combination of barriers that interact and reinforce one another within a system whose marginal value mechanism is already fundamentally different. Quantifying the gap at the operator level therefore requires a Norway-specific model rather than a transferred Nordic case study, and the institutional explanation of why the gap persists requires accounting for each of the four compounding conditions rather than isolating a dominant single cause. \utsett{autoref Section 2.5} formalities the two valuation perimeters that structure the quantification, and Chapter 3 specifies the model built to measure the gap they define.


%-------------------------------------------------------
%-------------------------------------------------------
%-------------------------------------------------------
\subsection{Commercial and Socio-Economic Value}

\utsett{Section opening statement 2.5}

%-------------------------------------------------------
\subsubsection{Defining the Two Perspectives}

The bedriftsøkonomisk valuation perimeter encompasses the cash flows that accrue to the district heating operator, and only those flows. On the revenue side, this perimeter includes spot-market arbitrage revenue from purchasing low-price electricity for power-to-heat conversion and displacing more expensive backup boiler dispatch during high-price hours; reserve-market payments for mFRR capacity and, where the asset qualifies technically, aFRR capacity, comprising both the capacity component paid regardless of activation and the activation payment received when the reserve is called; and avoided peak-load charges, to the extent that flexible dispatch reduces the operator's highest hourly power draw in a given billing period under the effektledd regime. On the cost side, the perimeter includes electricity purchased at the prevailing spot price, the effektledd network tariff, capital servicing on power-to-heat and storage assets, and variable operating costs including maintenance and staffing. The perimeter explicitly excludes the value that flexible operation creates for the grid operator, for other market participants, and for the broader energy system: these belong to the complementary perimeter. 

The socio-economic valuation perimeter encompasses all welfare effects attributable to district heating flexibility provision, regardless of which actor's balance sheet they appear on. In addition to the flows counted in the bedriftsøkonomisk perimeter, it includes avoided grid reinforcement costs, which accrue to the distribution system operator Elvia and ultimately to network tariff payers through a lower regulatory asset base; freed grid connection capacity that enables faster electrification of new loads, with welfare value captured by the connecting parties rather than the district heating operator; system price effects, whereby withdrawing power-to-heat demand from high-price hours clears the electricity market against a lower-cost marginal unit and redistributes consumer surplus across the bidding zone; displaced fossil or high-cost peaking generation and the corresponding reduction in $CO_2$ emissions; and, in the Norway-specific case, the conservation of reservoir water for cross-border export or intraday reserve provision, the value of which accrues primarily to hydropower producers and indirectly to export market participants rather than to the district heating operator \cite{askeland_balancing_2019}. The perimeter also incorporates an optionality component: the resilience value of thermal storage as a hedge against grid supply interruptions, a system reliability benefit with a diffuse beneficiary set. The conceptual basis for the socio-economic perimeter draws on the standard Norwegian framework for socio-economic analysis of public investments, adapted here to a private infrastructure operator context in which the externalities accrue to identifiable institutional actors rather than to anonymous social welfare.

\illu{FIGURE (§2.5.1): A monetary attributes schematic, companion to the physical system diagram in 2.4.1. Two side-by-side columns: bedriftsøkonomisk perimeter (operator ledger items) and samfunnsøkonomisk perimeter (all welfare items). Items that appear in the BØ column are shown in both; items exclusive to the SØ perimeter are marked distinctly. The visual should make immediately legible that SØ is a superset of BØ, and that the gap is the set of welfare items the operator cannot capture.}

These two perimeters are not two measurements of the same value. They are two accounting boundaries applied to the same physical asset and the same flexibility action. A single mFRR activation that reduces peak demand on the NO1 grid yields one value when assessed against the operator's income statement, and a larger value when assessed against total welfare. The difference is not a measurement error but a structural consequence of the institutional arrangements that determine which actor captures which cash flow. Anchoring the scale of that difference against prior work is difficult, because the field lacks a shared method. A systematic review of 123 district heating benefit and economic-assessment studies finds that flexibility, although identified as a benefit in fifteen of them, is monetized in none \cite{ahmed_district_2025}. The same review finds that welfare-level societal benefits are the least frequently assessed benefit category, and that no harmonized framework for valuing district heating benefits exists across the corpus. The choice of valuation perimeter is, consequently, not a presentational detail but the decision that fixes what the resulting figure means.

Two studies in the Nordic corpus apply paired perspectives to the same district heating asset and are the closest methodological precedents for the approach taken here. The first applies EnergyPLAN and energyPRO to the same Samsø asset and finds that the business-economically optimal storage size is only about 10 to 15 MWh, while up to roughly 3{,}200 MWh is beneficial from the island's societal perspective \cite{ostergaard_business_2019}. The authors describe this misalignment between the operator and societal optima as a gap to be overcome. The second combines a heat-merit-order business-economic analysis with a socioeconomic assessment of seasonal storage in Linz, and likewise finds that storage value depends on multi-function dispatch integration rather than on raw capacity \cite{moser_socioeconomic_2018}. Its societal layer rests partly on input-output multipliers rather than welfare externalities, however, so it transfers as a method rather than as a comparator. The Norwegian institutional conditions documented in \utsett{autoref: Sections 2.3 and 2.4}, the grid tariff structure, the energiloven §5-5 price cap, and the hydro-dominated value pathway) make direct transfer of either study's parameter values inadvisable, so the definitions applied here are constructed from the ground up for the Norwegian case. To the best of the author's knowledge, no paired bedriftsøkonomisk and samfunnsøkonomisk valuation of thermal flexibility has been conducted for Norwegian district heating at the firm level, a gap the methodology in \utsett{autoref: Chapter 3} is designed to fill.



%-------------------------------------------------------
\subsubsection{The Value Gap Mechanism}

In the resource-and-enabler terms of \utsett{ref: Section 2.1.1}, the value gap this thesis quantifies is what opens when a flexibility resource carries samfunnsøkonomisk value that no enabler converts into a bedriftsøkonomisk return. It operates through a four-step structural mechanism. First, thermal flexibility provision by a district heating operator reduces costs elsewhere in the energy system: avoided grid reinforcement at the distribution level, freed balancing capacity at the transmission level, and suppressed peak clearing prices that reduce consumer expenditure during high-demand events, each of which constitutes a real welfare gain attributable to the flexibility action \cite{mitridati_heat_2020}, \cite{mitridati_power_2018}, \cite{Thema_fjernvarme_2024}. Second, these cost reductions accrue primarily to Elvia as avoided distribution-grid capital expenditure, to Statnett as freed reserve capacity available for redeployment, and to electricity consumers as lower peak clearing prices. None of these actors are obligated under current market design to compensate the district heating operator for the service that generates them. This is the split-incentive configuration: societal-scale analysis shows that flexibility positively affects society, while industry does not capture a corresponding return \cite{sneum_barriers_2021}. The same divergence is documented directly for Norwegian network companies: a flexibility resource creates value for the transmission operator and wider society that the procuring company's own cost-benefit analysis does not capture \cite{saele_understanding_2023}. Third, the operator bears the full cost of flexible operation: electricity purchased at the prevailing spot price, effektledd peak charges that may rise rather than fall when a load-raising flexibility event sets a new monthly maximum, capital servicing on power-to-heat and storage assets, and the opportunity cost of heat not sold at the regulated price during arbitrage windows \cite{kirkerud_impacts_2016}, \cite{djorup_district_2020}. Fourth, no transfer mechanism currently bridges the gap between the samfunnsøkonomisk value generated and the bedriftsøkonomisk cost borne. The reduction of the mFRR bid floor to 1 MW Nordic-wide in March 2025 removed the bid-size exclusion that had kept most district heating plants out of the reserve market \cite{solberg_manedsrapport_nodate}. Yet prequalification cost, baseline and automation complexity, and the absence of an aggregation framework continue to deter sub-threshold participation \cite{monie_power--heat_2022}. No distribution-level product compensates peak-shaving, and district heating receives no nettnytte payment for the congestion it relieves \cite{Thema_fjernvarme_2024}. The effektledd structure does not exempt flexible loads from capacity charges during reserve activations, even when those activations reduce net system congestion. Finally, the price models available to the operator provide no route to pass flexibility value to customers or the wider system \cite{lygnerud_capturing_2024}. The result is a structural underinvestment equilibrium in which the operator's privately rational decision, to build out flexibility capacity and operate it actively, is attenuated by rules that assign the costs of flexibility to the actor who provides it while distributing its benefits to actors who bear none of the cost. This is the same divergence between operator-optimal and socially-optimal flexibility provision documented empirically on a single district heating asset, which is evidence that the gap is structural rather than hypothetical \cite{ostergaard_business_2019}. The equilibrium is not a market failure in the textbook sense of an externality that no institutional arrangement has yet recognized. It is a regulatory-design outcome that made sense under the conditions prevailing when the relevant rules were written, and that has not been updated to reflect the changed system conditions produced by the electrification transition, the growth of variable renewable generation in adjacent bidding zones, and the grid-constraint trajectory documented for the Norwegian system in \utsett{autoref: Section 2.4}. The gap is, in that sense, politically tractable: the numerical estimates in \utsett{autoref: Section 4.3} feed directly into the policy-instrument discussion in \utsett{ref: Section 5.3}.

The mechanism above describes the dominant direction of the gap: the samfunnsøkonomisk value of district heating flexibility exceeds the bedriftsøkonomisk value the operator can capture. This is the direction the Norwegian institutional arrangements predict and the pattern documented across the Nordic precedents. The directional claim warrants one qualification. Stylized Swedish district heating configurations with high power-to-heat penetration show that operator-optimized dispatch of stored power-to-heat can displace electricity co-generation from a competing biomass-CHP unit \cite{monie_power--heat_2022}. In that waste-to-energy configuration, biomass-CHP electricity output falls from $115 GWh$ to zero as power-to-heat capacity scales to 60 MW, with a corresponding loss of system balancing capacity that the operator's profit and loss does not register. The implication is not that the gap can run negative in the Norwegian case at current power-to-heat penetration, where Hafslund Celsio's co-generation comes from waste incineration dispatched as inflexible base-load rather than from a merit-order-displaceable biomass-CHP unit. Rather, at sufficiently high power-to-heat build-out, the assumption that samfunnsøkonomisk value always weakly exceeds bedriftsøkonomisk value requires re-examination. The analysis in \utsett{ref: Chapter 4} operates within the dominant-direction regime. The boundary condition under which the gap could reverse is flagged for the future work in \utsett{ref: Section 5.5}.


%-------------------------------------------------------
%-------------------------------------------------------
%-------------------------------------------------------
\subsection{Literature Review: Field Overview and Research Gap}

The preceding five sections have assembled the conceptual and empirical concepts this thesis utilizes: the vocabulary of flexibility \utsett{ref: 2.1}, the technology stack that supplies it \utsett{ref: 2.2}, the barrier taxonomy that constrains it \utsett{ref: 2.3}, the Norwegian system characteristics that make it a distinct case \utsett{ref: 2.4}, and the paired valuation perimeters that define what is being measured \utsett{ref: 2.5}. This closing section turns from synthesis to positioning. It is not a full systematic review. Its function is narrower: to make explicit how the source base was assembled, to map that base onto the dimensions that define the thesis's contribution, and to state precisely which region of that map remains unoccupied. The research-gap matrix in \utsett{ref: 2.6.2} carries the analytical weight of this positioning. The surrounding prose sets it up and reads it off, and the empty cell it identifies is the bibliographic form of the contribution claim made in \utsett{ref: 1.5}.

%---------------------------------------------------------
\subsubsection{How the Field Was Mapped}
The corpus underpinning this chapter was assembled through an informational rather than a structural review process. The starting point was a seed set of supervisor-supplied papers and the author's own keyword seeding. This was expanded in Research Rabbit and Consensus by snowballing across citation and co-citation links, and supplemented by targeted Scopus searches on the themes where the seed set was thin: grid tariff design, balancing-market access for demand-side assets, and Norwegian district heating specifically. 

This is an informational review, and is presented as such. Across 123 studies, the field has no harmonized method for monetizing district heating benefits \cite{ahmed_district_2025}, so a protocol-driven review would impose a comparability the underlying literature does not possess. A published systematic-methodology checklist is instead used as a coverage-hygiene instrument \cite{atabaki_systematic_2025}, to confirm that the assembled set spans the temporal, spatial, economic, uncertainty, and sector-coupling dimensions a flexibility valuation must address, rather than as the review method itself. The assembled corpus is, by composition, overwhelmingly Danish- and Nordic-led. The Norwegian contributions are largely case-applications of frameworks developed for the Danish and wider Nordic systems rather than independent methodological work. Even at the field level, the intersection of district heating with sector-coupling policy is thin, explicit in only one of the 25 studies in a recent sector-coupling scoping review \cite{vanhanen_energy_2024}. The research-gap matrix encodes this geographic skew directly.




%---------------------------------------------------------
\subsubsection{Research-Gap Matrix}

The corpus can be positioned on two dimensions that jointly define the thesis's contribution. The first is analytical scope: whether a study models a single operator or plant, an aggregate market or region (system-level), or compares multiple systems or reviews the field (multi-system). The second is valuation perspective, the distinction formalized in \utsett{ref section 2.5:} whether a study measures only the (bedriftsøkonomisk) perspective, only (samfunnsøkonomisk), or both on consistent scenarios (paired). Geographic specificity, the third dimension relevant to the contribution claim, is encoded by color rather than as a third axis, so that the Norwegian work is visually separable from the broader Nordic and European corpus. Figure \utsett{autoref gap matrix} plots the principal sources on this grid.

\todo{Research-gap matrix for the valuation of thermal flexibility in district heating. Analytical scope (columns) against valuation perspective (rows); marker color encodes geography (Norwegian, other Nordic, EU/global). Among prior work, the firm-level $\times$ paired cell holds a single marker, and it is Danish (Østergaard et al. 2019); no Norwegian study appears anywhere in the paired row, which is the cell this thesis enters. Field-spanning reviews populate the multi-system column; purely conceptual or baseline sources that take no bedriftsøkonomisk or samfunnsøkonomisk valuation stance (Sneum 2021; Werner 2017; Fernqvist 2023; Ma 2020; Selvakkumaran 2021) are part of the corpus but are not plotted as valuation markers. Source assignments are indicative and follow the project's source index.}

\illu{BUILD Fig. literature matrix Spec to reproduce: axes: X = analytical scope [firm-level | system-level | multi-system/review], Y = valuation perspective [paired BØ+SØ | SØ-only | BØ-only]; geography by colour (NO = dark red, Nordic = blue, EU/global = grey). Cells: (firm x paired) THIS THESIS (NO, highlighted focal cell) + Østergaard 2019 (DK, singular paired paper); (firm x SØ-only) Grønli/AFRY 2025 (NO), Moser 2018 (AT/EU -- classified SØ-only here, not paired, because its SØ side rests on input-output multipliers rather than a welfare CBA); (firm x BØ-only) Kirkerud 2016 (NO), Trømborg 2017 (NO), Johannsen 2021 (DK), Javanshir 2022, Monie 2022, Tan 2022; (system x SØ-only) THEMA 2024, Vista 2023, Askeland 2019, SINTEF ZEN-47, Sandberg 2019, Bergaentzlé 2022, Mitridati 2020, Mitridati \& Taylor 2018 (formal virtual-storage equivalence), Kirkerud 2017, Atabaki 2026 (SE), Gea-Bermúdez 2021 (EU), Boldrini 2022 (EU, DH balancing-market potential), NOU 2023:3 / Mer av alt (NO national), OED 2022 / Har vi fleksibilitet nok (NO, independent of THEMA 2024), RME 2020 (NO regulator critique), Hofmann 2025 (Oslo); (system x BØ-only) Tan 2022 (borderline); (multi-system x SØ-only) Vanhanen 2024, NVE LA23 2023, Ahmed 2025, Atabaki 2025 (categorisation matrix); (multi-system x BØ-only review) Golmohamadi 2022; other cells empty. open: geography x perspective, or a compact 2x2 of single-vs-paired against NO-vs-other -- lock the structure before finalising. Verify every plotted source against the project's source index.}

The matrix makes the gap legible without recourse to the prose, and two field-wide observations frame it. First, a recent critical review of local flexibility markets in Europe notes that "the future electrification of industry and the transformation of district heating systems remain comparatively underexplored, despite their clear relevance as new sources of flexibility within a sector-coupling perspective" \cite{scrocca_local_2026}. The gap the matrix identifies is therefore recognized at the European level, not only in the Norwegian-specific corpus. Second, the gap is methodological as much as geographic. Across 123 district heating benefit studies, flexibility is quantified in only four and monetized in none \cite{ahmed_district_2025}, so the paired row is sparse across the whole field, not merely for Norway. Within that field, the two single-perspective rows are both populated for Norway. The bedriftsøkonomisk row has firm-level Norwegian precedent \cite{kirkerud_impacts_2016}, and the samfunnsøkonomisk row carries the system-level Norwegian welfare diagnosis of the THEMA report \cite{Thema_fjernvarme_2024}, with further Norwegian system-level studies populating the same row in Figure \utsett{ref gap matrix}. The paired row, by contrast, is almost empty. The only study in the corpus that sets a bedriftsøkonomisk and a samfunnsøkonomisk welfare assessment of district heating flexibility on consistent scenarios sits in the firm-level $\times$ paired cell, as a single Danish marker on a small biomass island with no hydropower interaction \cite{ostergaard_business_2019}. A separate study pairs a plant-level monetization with a regional macroeconomic layer for an Austrian storage case, but that layer rests on input-output multipliers rather than a welfare cost-benefit analysis, so the matrix places it in the samfunnsøkonomisk row rather than the paired one \cite{moser_socioeconomic_2018}. To the extent the cell that defines this thesis is occupied by prior work at all, it is occupied by a single study that is methodologically reusable but substantively non-transferable, and no Norwegian marker appears anywhere in the paired row.


%---------------------------------------------------------
\subsubsection{Where This Thesis Fits and What It Adds}
This thesis occupies the firm-level $\times$ paired cell of \utsett{ref gap-matrix}, applied to the Norwegian district heating case: a single operator's flexibility assessed against both the bedriftsøkonomisk and the samfunnsøkonomisk perimeter on consistent scenarios. It draws on the two nearest populated regions of the map. From the single-perspective Norwegian work it takes the operator-ledger modeling precedent \cite{kirkerud_impacts_2016} and the welfare-side institutional diagnosis of the THEMA report \cite{Thema_fjernvarme_2024}. From the paired Nordic work it takes the methodological template of the Danish Samsø study \cite{ostergaard_business_2019}, adopted for its two-perimeter structure rather than for its biomass-island results. The precise claim the matrix supports is therefore not that the value of thermal flexibility is unstudied, but that, to the best of the author's knowledge, no paired bedriftsøkonomisk and samfunnsøkonomisk valuation of thermal flexibility has been conducted for Norwegian district heating at the firm level. What makes the cell more than a geographic blank is the regulatory detail that defines it: although generic grid-tariff effects on power-to-heat are well studied, the Norwegian effektledd capacity tariff central to the operator-side cost in §2.3.2 has been characterized in policy and operator-facing analysis rather than analyzed as a flexibility barrier in its own right. The thesis also quantifies before it monetizes: a technical flexibility envelope (Section 3.2) sizes the physical resource in MWh and MW and supplies the common denominator for the paired valuation, directly addressing the field-wide finding that flexibility has been monetized in none of the 123 reviewed studies \cite{ahmed_district_2025}. The methodology developed in Chapter 3 is constructed specifically to fill that cell.

\avklar{The novelty claim; Scopus + Google Scholar + Oria check with the exact query phrasing of the cell description.}


%---------------------------------------------------------
\subsubsection{Adjacent Gaps Left for Future Work}
Naming the cell the thesis fills also sharpens the adjacent cells it does not. Three are worth flagging, both to anchor the future-work discussion in §5.5 and to keep the contribution claim precise.

\begin{itemize}
  \item \textbf{Multi-system paired valuation.} The thesis takes Oslo as its primary case; extending the paired BØ/SØ assessment across Trondheim, Bergen, and smaller municipal systems is a distinct research area, not a marginal extension, because the storage and capacity asymmetries documented in sections 2.3.1 and §2.4.1 break the single-case parameterization.
  \item \textbf{End-user welfare and distributional effects.} The welfare consequences of flexibility activation for the heat customer fall outside the scope set in §1.4, and remain almost entirely unmonetized in the field: end-user and distributional benefits are valued in only a handful of the 123 studies in the systematic review \cite{District Heating Benefits and Economic Assessment Methods: A Systematic Review and the Role of Emerging Technologies - Ahmed - 2025}.
  \item \textbf{Business-model and instrument design.} This thesis names the missing transfer mechanism (§2.5.2) but does not design it. The concrete shape of a capacity-payment product, an aggregation contract, or a bilateral DSO flexibility agreement remains open. The closest precedent for that design question is a systematic review of district-energy business models \cite{selvakkumaran_how_2021}.
\end{itemize}









\newpage
